\section{Bounded changes of measure and the generator domain}
\label{sec:drift}

Our first task is to control the weighted generator before assigning
any curvature bound to the weighted space. A bounded Lipschitz weight
preserves weak gradients and the generator domain. The less immediate
point is that bounded functions with bounded weighted Laplacian are
Lipschitz; the resolvent argument below supplies this regularity from
the original heat flow.

\begin{lemma}[Bounded drift]
\label{weight:resolvent}
Let \((M,d,m)\) be a full-support \(\RCD(K,N)\) space, with \(1<N<\infty\), and
let \(V\) be bounded and globally Lipschitz. Set
\(m_V=e^V m\), and let \(\Delta_V\) be the generator of the
Cheeger energy for \(m_V\), with the nonpositive sign convention.
The Sobolev spaces and their weak gradients agree, and
\begin{equation}
\label{weight:generator}
D(\Delta_V)=D(\Delta),\qquad
\Delta_Vu=\Delta u+\Gamma(V,u).
\end{equation}
The two generator graph norms are equivalent. Moreover,
\begin{equation}
\label{weight:multiplier-class}
\xi\in D(\Delta_V),\quad
\xi,\Delta_V\xi\in L^\infty
\quad\Longrightarrow\quad
\xi\text{ has a globally Lipschitz representative}.
\end{equation}
\end{lemma}

\begin{proof}
Since \(e^V\) is bounded above and below by positive constants,
the classes of test plans defining weak upper gradients agree:
a compression bound for one measure gives a compression bound
for the other. Thus the weak gradients agree, and the two
\(W^{1,2}\) norms are equivalent. In particular, the weighted
Cheeger energy is quadratic. Multiplication by \(e^V\) or
\(e^{-V}\) preserves \(W^{1,2}\). In the weak definition of the
generator, replacing a test function \(\eta\) by \(e^V\eta\)
gives \eqref{weight:generator} in one direction. Replacing it by
\(e^{-V}\eta\) gives the converse. All drift terms are in \(L^2\)
because \(|\nabla V|\) is bounded.
The energy identity
\[
\|\nabla u\|_{L^2(m)}^2=-\int u\Delta u\dd m
\]
and its weighted counterpart bound the gradient by the respective
graph norm. Together with \eqref{weight:generator}, this proves
graph-norm equivalence.

Here is a resolvent proof of the last assertion, before imposing
any curvature condition on the weighted space. Let \(P_t\) be
the original heat semigroup, and put
\[
R_a h=\int_0^\infty e^{-at}P_t h\dd t=(a-\Delta)^{-1}h.
\]
For \(h\in L^2\cap L^\infty\), spectral calculus, the Markov
property, and the heat gradient estimate give, for all sufficiently
large \(a\),
\begin{align}
\label{weight:resolvent-bounds}
\|R_a h\|_p&\le a^{-1}\|h\|_p\quad(p=2,\infty),\\
\|\nabla R_a h\|_2&\le C a^{-1/2}\|h\|_2,
&\|\nabla R_a h\|_\infty&\le C_Ka^{-1/2}\|h\|_\infty.
\nonumber
\end{align}
For the last estimate one integrates the reverse-Poincar\'e
bound \(\|\nabla P_t h\|_\infty\le
C_K(t^{-1/2}+1)e^{C_Kt}\|h\|_\infty\).
These are estimates for the original \(\RCD\) heat semigroup.

Consider the Banach space
\[
\mathcal E=\{z\in W^{1,2}:z,|\nabla z|\in L^\infty\},\qquad
\|z\|_{\mathcal E}=\|z\|_2+\|\nabla z\|_2
                         +\|z\|_\infty+\|\nabla z\|_\infty.
\]
Completeness follows from closedness of the weak differential
and completeness of \(L^\infty\). Given \(h\in L^2\cap L^\infty\),
the map
\[
z\longmapsto R_a\bigl(h+\Gamma(V,z)\bigr)
\]
is a contraction on \(\mathcal E\) when \(a\) is sufficiently
large, by \eqref{weight:resolvent-bounds} and
\(\|\Gamma(V,z)\|_p\le\|\nabla V\|_\infty\|\nabla z\|_p\).
Its fixed point lies in \(D(\Delta)\) and solves
\((a-\Delta_V)z=h\). The weighted energy identity gives
uniqueness of this solution in \(D(\Delta_V)\).
For \(\xi\) as in \eqref{weight:multiplier-class}, take
\(h=(a-\Delta_V)\xi\). It belongs to \(L^2\cap L^\infty\),
so uniqueness identifies \(\xi\) with the fixed point.
The original Sobolev-to-Lipschitz property now proves the claim.
\end{proof}

